% Material migrated from the former appendix into the main body for IPDPS.
% All added prose is review-marked so coauthors can inspect it before finalization.

\anurag{Appendix - A Background: AllReduce, AllGather, and Point-to-Point communication}

\review{\noindent\textbf{Communication patterns under the three parallelism strategies.}
Tensor, pipeline, and data parallelism expose different communication structures, and these differences determine both the amount of data moved and where communication energy appears in an inference run. In tensor parallelism, the main collective is ring AllReduce. GPUs are arranged in a logical ring and each GPU exchanges data only with its immediate neighbors, avoiding an all-to-all exchange. Ring AllReduce proceeds in two phases. In the \emph{ReduceScatter} phase, each GPU divides its tensor into $p$ chunks, sends one chunk to its successor, receives one from its predecessor, and reduces the received values with its local chunk. After $p-1$ steps, every GPU holds one fully reduced shard. In the subsequent \emph{AllGather} phase, those reduced shards are circulated for another $p-1$ steps until every GPU reconstructs the complete reduced tensor. Thus, tensor parallelism repeatedly combines data movement with synchronization, which makes its communication energy particularly sensitive to GPU count, message size, and rank-to-rank waiting behavior.}

\review{Pipeline parallelism instead communicates activations between adjacent stages. After stage $i$ completes the forward computation for a microbatch, it sends its output activations to stage $i+1$ through a point-to-point device transfer, e.g., over PCIe or NVLink. During inference there are no gradient exchanges, so these hop-local activation sends and receives are the primary communication events. The transfer lies directly on the dependency path between the last forward kernel of stage $i$ and the first kernel of stage $i+1$, making the stage boundary the natural unit for communication-energy accounting.}

\review{Data parallelism has a different structure because each replica executes a complete forward pass on its local input independently. In our inference setup, communication occurs primarily at the end of the forward pass, where a single AllGather collates the final outputs (e.g., logits or token scores) across replicas. The output exchange begins after the local output layer completes and the concatenated result is then handed to the host for downstream sampling or post-processing. Compared with tensor and pipeline parallelism, this tail-end communication is less frequent and more structured, which is one reason data-parallel energy is easier to attribute.}
